\documentclass[11pt]{article}
\usepackage[a4paper,margin=28mm]{geometry}
\usepackage{amsmath,amssymb,amsthm,mathtools,booktabs,hyperref}
\newtheorem{theorem}{Theorem}[section]
\newtheorem{lemma}[theorem]{Lemma}
\newtheorem{proposition}[theorem]{Proposition}
\newtheorem{corollary}[theorem]{Corollary}
\theoremstyle{remark}\newtheorem{remark}[theorem]{Remark}
\newcommand{\Q}{\mathbb Q}\newcommand{\Z}{\mathbb Z}
\newcommand{\R}{\mathbb R}\newcommand{\vpp}{v_p}
\title{A uniformly bounded real normalization for the scalar elliptic detector}
\author{}\date{}
\begin{document}\maketitle
\begin{abstract}
We give a scalar detector whose real preparation is uniform in the
elliptic point carrying the large index.  The construction replaces an
unbounded formal parameter by a bounded rational function and uses a
positive rational normalization bounded away from zero.  Its local
valuation implications are unchanged outside an explicitly defined
finite set of primes.  This removes a real approximation requirement
from the scalar certificate.  It does not assert a reduction in the
number of variables: prescribed indices must still satisfy the
independent reduction-order conditions at the finite contact primes.
\end{abstract}

\section{The elliptic curve and a bounded parameter}

Put $K=\Q(\sqrt2)$, $D=(3-\sqrt2)^2$, and
\[
 E:\quad y^2=x^3-Dx.
\]
Write $\sigma$ for the nontrivial automorphism of $K/\Q$.
An embedding of $K$ into $\R$ will be denoted by $\iota$.
The identity component of $E_\iota(\R)$ consists of the point at
infinity and the affine points with $x\ge\sqrt{\iota(D)}$.
In particular, every nonzero double has such an abscissa.
For an affine non-torsion point define
\[
 f(Q)=\frac1{x(Q)},\qquad z(Q)=-\frac{x(Q)}{y(Q)},\qquad
 w(Q)=\frac{y(Q)}{x(Q)^2+D}.
\]

\begin{lemma}\label{lem:bounded}
At either real embedding and every nonzero double $Q$,
\[
 0<\iota(f(Q))\le\frac1{\sqrt{\iota(D)}}<1,
 \qquad |\iota(w(Q))|<1.
\]
The denominator $x(Q)^2+D$ never vanishes at a $K$-point.
Moreover,
\begin{equation}\label{eq:parameter-unit}
 \frac{w(Q)}{z(Q)}
 =-\frac{1-Df(Q)^2}{1+Df(Q)^2}.
\end{equation}
At a good nonarchimedean place where $Q$ lies in the formal group,
the right side is a unit and is congruent to $-1$ in the residue field.
\end{lemma}
\begin{proof}
The assertion about $f$ follows from the description of the real
identity component and $3-\sqrt2>1$.  For $x>\sqrt D$, evaluated
at either embedding,
\[
 w^2=\frac{x(x^2-D)}{(x^2+D)^2}<\frac1x
 \le\frac1{\sqrt D}<1.
\]
The same bound holds at the endpoint, where $w=0$.
Total positivity of $D$ proves nonvanishing of $x^2+D$.
Using $y^2=x(x^2-D)$ gives~\eqref{eq:parameter-unit}.
At a good formal point, $v(f)>0$ and $D$ is integral, so the
numerator and denominator on the right are units.
\end{proof}

\begin{corollary}\label{cor:ratio}
Suppose $P$ is a fixed non-torsion point and both $[c]P$ and
$[h]P$ are in the formal group at a good place, with $c,h$ nonzero
integers.  Then
\[
 v\!\left(\frac{w([c]P)}{w([h]P)}\right)
 =v\!\left(\frac{z([c]P)}{z([h]P)}\right)=v(c/h).
\]
\end{corollary}
\begin{proof}
The first equality follows from Lemma~\ref{lem:bounded}.
For the second, the formal logarithm divided by $z$ is a unit
on the sufficiently deep formal neighborhood used in the scalar
detector, and the formal logarithm is additive.  At the large
contact primes used here this is the usual first formal neighborhood.
Equivalently, the two logarithms are nonzero integer multiples
of the logarithm of a common nonzero formal multiple of $P$.
\end{proof}

\section{A positive normalization bounded away from zero}

Define the integer polynomials
\[
 \begin{split}
  \beta(X)&=216X^4+450459,\\
  \rho(X)&=\beta(X)(\beta(X)-1)(\beta(X)-2401)\\
   &\hspace{21mm}\cdot(32\beta(X)+343)(2\beta(X)-1).
 \end{split}
\]
Put
\[
 \ell=76249,\quad q=1021,\quad a_0=8\cdot108^5,\quad
 A_0=66\cdot450459^5.
\]
The letters $\ell$ and $q$ denote two fixed primes.  Both are
avoidance primes for $\rho$: the accompanying exact certificate
checks $\rho(x)\ne0$ at every residue class modulo either prime.

\subsection{A fixed polynomial denominator}

Define
\begin{equation}\label{eq:h0}
 h_0(X)=2239488X^{10}+11663498286X^6
              +\frac{9330368884708461}{512}X^2.
\end{equation}
Direct expansion gives
\begin{equation}\label{eq:F}
 \begin{split}
 F(X):=\rho(X)-6h_0(X)^2
  ={}&\frac{10882407499318210844052639}{64}X^8\\
    &+\frac{110270372953940856744489540803781}{131072}X^4\\
    &+1180712527837405404989927078052.
 \end{split}
\end{equation}
Let $N$ be the sum of these three positive coefficients.  Let $T$
be the smallest integer with
\[
 T\ge\max(\lceil N\rceil+1,2239488),\qquad T\equiv2\pmod\ell.
\]
Explicitly,
\[
                 T=1181553994055396276117553880104.
\]
Set
\begin{equation}\label{eq:fixed-denominator}
 \begin{gathered}
 h=h_0+T,\qquad H=512h\in\Z[X],\qquad
 \delta=\rho-6h^2\in\Z[1/2][X],\
 J=\operatorname{Res}_X(\rho,H),\qquad
 D_0=\left\lceil\sum_i|[X^i]\delta|\right\rceil,\qquad
 \mu=\frac{3\cdot2239488}{8}.
 \end{gathered}
\end{equation}
These quantities are fixed; none depends on the later exponent $L$.

\begin{lemma}\label{lem:denominator}
The integer $J$ is nonzero.  For every real $X$,
\[
 \delta(X)<0,\qquad
 \mu(1+X^2)^5<|\delta(X)|\le D_0(1+X^2)^5,
\]
and
\begin{equation}\label{eq:rho-ratio}
 \frac{a_0}{D_0^2}
 <\frac{\rho(X)}{\delta(X)^2}<\frac{A_0}{\mu^2}.
\end{equation}
The reduction of $\delta$ has no root modulo $\ell$ and its
degree-ten leading coefficient is an $\ell$-unit.
\end{lemma}
\begin{proof}
Put $\psi=1+X^2$.  As $X^4+1\ge\psi^2/2$,
\[
 108\psi^2\le\beta(X)\le450459\psi^2,
 \qquad a_0\psi^{10}<\rho(X)<A_0\psi^{10}.
\]
All coefficients of $h_0$ and $F$ are nonnegative, and
$F(X)\le N\psi^4$.  Hence
\[
 0<\sqrt{\rho(X)/6}-h_0(X)
 =\frac{F(X)}{6(\sqrt{\rho(X)/6}+h_0(X))}
 \le\frac{N}{\sqrt{6a_0}\,\psi}<N.
\]
It follows that $h-\sqrt{\rho/6}>T-N\ge1$.  Thus
\[
 |\delta|=6(h-\sqrt{\rho/6})(h+\sqrt{\rho/6})>6h.
\]
The inequalities $T\ge2239488$ and
$(1+X^2)^5\le16(1+X^{10})$ imply $6h\ge\mu\psi^5$.
The coefficient norm of the degree-ten polynomial $\delta$ gives
the upper bound for $|\delta|$.  These facts prove~\eqref{eq:rho-ratio}.

If $\rho(\alpha)=0$, then $\alpha^4$ is a negative rational
number, so $\alpha^2$ is purely imaginary.  Formula~\eqref{eq:h0}
shows that $h_0(\alpha)$ is purely imaginary.  Since $T>0$,
$H(\alpha)\ne0$.  Therefore $J\ne0$.
Finally the two assertions modulo $\ell$ are checked on all
$76249$ residue classes in the supplied certificate, using
$T\equiv2\pmod\ell$.  The leading coefficient is
$-12\cdot2239488\,T$, which is an $\ell$-unit.
\end{proof}

\begin{lemma}\label{lem:odd-valuations}
For every integer $X$ and every prime $p\ne2,3$, if
$v_p(\rho(X))$ is odd, then
\[
                    v_p(\delta(X))\le2v_p(J).
\]
Consequently $J^4\rho(X)/\delta(X)^2$ has no negative odd
valuation at any prime $p\ne2,3$.
\end{lemma}
\begin{proof}
The polynomials $\rho,H$ are integral.  The resultant identity
implies
\[
 \min\{v_p(\rho(X)),v_p(H(X))\}\le v_p(J).
\]
For $p\ne2,3$ the denominator $512$ and the coefficient $6$ are
units.  If $u=v_p(\rho(X))$ is odd and $v=v_p(h(X))$, the two
terms in $\delta=\rho-6h^2$ have different valuation parity.
Thus $v_p(\delta)=\min(u,2v)$.
If $u\le v_p(J)$ the desired bound is immediate.  Otherwise
$v\le v_p(J)$, which gives the same bound.  It follows that
\[
 v_p(J^4\rho/\delta^2)\ge4v_p(J)+u-4v_p(J)=u>0.
\]
If $u$ is even, the valuation of the displayed rational number
is even, which proves the last assertion.
\end{proof}

\subsection{The normalization and its local properties}

Write
\[
 U_0=\frac{A_0}{\mu^2},\qquad
 c_*=102^6J^4,\qquad
 R_{\min}=\frac{a_0}{D_0^2q^2U_0}.
\]
For each odd $L\ge3$, let $k(L)$ be the smallest nonnegative
integer satisfying
\[
                  q^{2k(L)}>c_*\ell^L U_0,
\]
and put
\begin{equation}\label{eq:normalizer}
 R_L(X)=\frac{c_*\ell^L\rho(X)}{q^{2k(L)}\delta(X)^2}.
\end{equation}

\begin{proposition}\label{prop:normalizer}
For every real $X$ and every odd $L\ge3$,
\[
                    0<R_{\min}<R_L(X)<1.
\]
The constant $R_{\min}$ is independent of $L$.  For every rational $X$,
\begin{equation}\label{eq:ell-depth}
                   v_\ell(R_L(X))=L+4v_\ell(J).
\end{equation}
For every integer $X$, the reduced denominator of $R_L(X)$ is a
square, $R_L(X)$ is a square in each of $\Q_2,\Q_3,\Q_{17}$, and
\[
 v_2(R_L)\ge4,\qquad v_3(R_L)\ge4,\qquad v_{17}(R_L)\ge4.
\]
If $p\nmid 2\cdot3\cdot17\cdot q\ell J$ is a positive contact
prime of $\rho(X)$, then $v_p(R_L(X))=v_p(\rho(X))$.
\end{proposition}
\begin{proof}
Since $U_0>1$ and $c_*\ell^L>1$, one has $k(L)\ge1$.
Its minimality gives
\[
 \frac1{q^2U_0}\le\frac{c_*\ell^L}{q^{2k(L)}}<\frac1{U_0}.
\]
Combine this with~\eqref{eq:rho-ratio} for the real assertions.
At an $\ell$-integral rational $X$, both $\rho(X)$ and $\delta(X)$
are units.  If $v_\ell(X)<0$, their unit leading coefficients give
$v_\ell(\rho(X))=20v_\ell(X)$ and
$v_\ell(\delta(X))=10v_\ell(X)$.  Both cases yield
\eqref{eq:ell-depth}.

For integral $X$, direct residue calculations give
\[
 \rho(X)\in\Q_2^{\times2}\cap\Q_3^{\times2}\cap\Q_{17}^{\times2},
 \quad v_2(\rho)=v_3(\rho)=2,\quad v_{17}(\rho)\in\{0,2\}.
\]
At $2$ and $3$ the valuation of $6h(X)^2$ is odd, whereas the
valuation of $\rho(X)$ is even.  Thus $v_p(\delta)\le v_p(\rho)=2$.
At $17$, the number $6$ is a nonsquare unit; since $\rho(X)$ is a
nonzero local square, cancellation is again impossible in
$\rho-6h^2$, and $v_{17}(\delta)\le v_{17}(\rho)\le2$.
The factor $102^6$ therefore supplies the stated local depths.
The number $\ell$ is a local square at $2,3,17$, and all remaining
factors in~\eqref{eq:normalizer} are rational squares, so the local
squareclasses are as claimed.

At $p\ne2,3,q,\ell$ apply Lemma~\ref{lem:odd-valuations}; the
additional factors are squares and do not lower the valuation.
At $q$, $\rho(X)$ is a unit by the avoidance certificate, so
$v_q(R_L)$ is even.  At $\ell$ it is positive by~\eqref{eq:ell-depth}.
At $2,3$ it is even and positive, as already proved.  Thus no
negative odd valuation occurs, proving the square-denominator
assertion.  Finally, at a positive contact prime outside the
displayed fixed set, the resultant identity makes $H(X)$ a unit.
The relation $\delta=\rho-6h^2$ then makes $\delta(X)$ a unit,
which proves the last assertion.
\end{proof}

\section{Uniform real preparation of the norm tree}

Let $\mathcal U$ be a fixed neighborhood in the compact real closure
of the relevant multiples of $P$, chosen so that
\begin{equation}\label{eq:H-open}
 |w(H)|\ge\eta,\qquad |\sigma w(H)|\ge\eta
 \quad(H\in\mathcal U)
\end{equation}
for some fixed $0<\eta\le1$.
Such a neighborhood exists around any non-torsion multiple in the
identity component of that closure.  It is compatible with arbitrary
fixed congruence conditions that select a finite-index subgroup:
the closure of every such subgroup contains the identity component
of the original compact closure.

Let $C_S=S-P$ be any nonzero double and $H\in\mathcal U$ a nonzero
double.  Write the rational components, in the basis $1,\sqrt2$, of
\[
             f(H),\qquad f(C_S),\qquad w(C_S)/w(H)
\]
as $g_0,\ldots,g_5$.  Set $m_0=\cdots=m_3=1$ and
$m_4=m_5=20$.  Fix an integer $K$ such that
\begin{equation}\label{eq:K}
 \ell\nmid K,\qquad
 K\ge\frac{4}{\eta\min(R_{\min},R_{\min}^{20})}.
\end{equation}
For $R=R_L(X)$ define
\[
 u_i=\frac{g_i}{K R^{m_i}},\qquad N_R(u,v)=u^2+Rv^2,
\]
and the same norm tree as in the scalar detector:
\[
 \begin{gathered}
 w_0=N_R(1,u_0),\quad w_1=N_R(u_1,u_2),\quad
 w_2=u_3,\quad w_3=N_R(u_4,u_5),\
 W=N_R\bigl(N_R(w_0,w_1),N_R(w_2,w_3)\bigr).
 \end{gathered}
\]
The constants $K$ and $R_{\min}$ are independent of $L$.

\begin{theorem}\label{thm:uniform}
For every real $X$ and every pair $C_S,H$ satisfying the preceding
conditions,
\[
 1\le W\le\frac{86249}{65536}<\frac32,
 \qquad 0<R_L(X)W^2<\frac94<17.
\]
These conclusions impose no real approximation condition on $C_S$.
At every contact prime outside the fixed set of prime divisors of
$2\cdot3\cdot7\cdot17\cdot q\ell KJ$, the local valuation implications
of the scalar detector are unchanged by replacing $z$ with $w$
and dividing its six component tests by $K$.
At the fixed prime $\ell$, the detector gives the index depth
\[
                  v_\ell(c)\ge19L+1,
       \qquad C_S=[c]P.
\]
\end{theorem}
\begin{proof}
Lemma~\ref{lem:bounded} and~\eqref{eq:H-open} bound the absolute
value of each $g_i$ by $1/\eta$: for a number $u+v\sqrt2$,
its rational components are half the sum and $1/(2\sqrt2)$
times the difference of its two real conjugates.
The choice of $K$ therefore gives $|u_i|\le1/4$.
Since $0<R<1$,
\[
 \begin{split}
 1\le w_0\le17/16,\quad |w_1|\le1/8,\quad
 |w_2|\le1/4,\quad |w_3|\le1/8,\\
 1\le N_R(w_0,w_1)\le293/256,\qquad
 0\le N_R(w_2,w_3)\le20/256.
 \end{split}
\]
Squaring these last bounds proves the asserted bound for $W$.

At a good consumed contact prime, $K$ and $\delta(X)$ are units;
the latter assertion follows from the resultant identity whenever
the input is integral at that prime.  Thus the contact valuation
of $R$ agrees with the prescribed contact valuation of $\rho$.  At such a contact prime, the positive contact of $\beta(X)$
also implies that $X$ is integral, since $p\nmid216$.
The order-one component tests put
$H,C_S$ in the formal groups exactly as before.  Lemma~\ref{lem:bounded}
and Corollary~\ref{cor:ratio} then show that the two order-twenty
tests have precisely the required ratio valuation.

More explicitly, if $\nu=v_\ell(R)>0$ is odd, a node
$u^2+Rv^2$ has valuation equal to the smaller term valuation,
because the two terms have different parity.  The pole obstruction
for $RW^2$ gives $v_\ell(W)>-\nu/2$.  Following every path through
the norm tree shows that each normalized leaf has valuation
greater than $-\nu$.  Since $\ell\nmid K$, both components of
$w(C_S)/w(H)$ have valuation greater than $19\nu$.  The exact
formal ratio consequently gives $v_\ell(c/h)\ge19\nu+1$.
Since $h\in\Z$, Proposition~\ref{prop:normalizer} yields the
claimed bound.
\end{proof}

\begin{remark}[Finite exceptions and the contact-height constant]
Suppose the contact-height theorem is written as
$h(\beta)\le H M(\beta)^{19}$, where $M(\beta)$ is the squarefree
product of the consumed contact primes.  Removing a fixed finite
set $F$ changes the bound only to
\[
 h(\beta)\le H\left(\prod_{p\in F}p\right)^{19}
                       M_F(\beta)^{19}.
\]
Thus the finite exceptional primes of this normalization change a
fixed constant and the resulting sufficient index threshold.
They do not require an additional variable or an additional
arithmetic existence theorem.  No claim that the displayed constant
or coefficients are small is intended.
\end{remark}

\section{The remaining finite-place condition}

The real estimate does not replace the finite-place choice of the
index $c$.  If $P$ has good reduction at a place $\mathfrak p$ and
its reduction has order $d_{\mathfrak p}$, then
\[
          [c]P\text{ belongs to the formal group at }\mathfrak p
          \quad\Longleftrightarrow\quad d_{\mathfrak p}\mid c.
\]
The integer $d_{\mathfrak p}$ can have prime factors different
from the residue characteristic.  Consequently $p^e\mid c$ does
not by itself imply formal-group membership.  In particular,
a proposed prescribed index
\[
                       c=M\rho(X)^e
\]
must still be shown to be divisible by the reduction order at
every variable contact prime.  The uniform real normalization
proved here leaves this arithmetic requirement explicit.

\section*{Verification}

The accompanying standard-library Python program
\texttt{check\_uniform\_scalar\_normalization.py} verifies the finite
residue assertions, the exact norm-tree bound, and the constant
inequalities.  The resultant is evaluated independently by a Sylvester determinant
and a product of five quartic-factor formulas.  The displayed proofs,
rather than finite sampling,
establish the assertions for all real or rational inputs.

\begin{thebibliography}{9}
\bibitem{Source}
\emph{A non-torsion construction in twenty-seven variables},
project paper, version 12, 9 November 2026.
\bibitem{Explicit}
\emph{An explicit non-torsion construction in twenty-nine variables},
project paper, version 12, 9 November 2026.
\end{thebibliography}
\end{document}
